\chapter{Pendahuluan}
\section{Latar Belakang}
\par Termodinamika kuantum merupakan ekstensi dari termodinamika klasik dan fisika statistik untuk menjelaskan perilaku termal pada skala sangat kecil di mana hukum-hukum fisika klasik runtuh \cite{Deffner2019Thermo}. Salah satu aplikasi dari termodinamika kuantum adalah mesin panas kuantum. Sebuah sistem siklik yang bertujuan untuk mengubah panas menjadi kerja menggunakan partikel kuantum sebagai zat kerjanya. Mesin panas kuantum memiliki beberapa properti yang unik. Sebagai contoh, sebuah siklus mesin panas kuantum mampu untuk memperoleh efisiensi yang lebih dari siklus Carnot klasik yang diketahui sebagai siklus mesin panas klasik paling efisien \cite{Kosloff2014Engines}. Properti unik ini disebabkan oleh hukum-hukum termodinamika yang berlaku secara berbeda untuk partikel kuantum. Sehingga, untuk mengkonstruksi sebuah mesin panas kuantum, perlu untuk dilakukan penyetaraan hukum-hukum termodinamika klasik agar bisa bekerja pada sistem kuantum. Dari hukum yang dimodifikasi tersebut, bisa dinyatakan proses-proses termodinamika seperti isotermal, isovolume, isobarik, dan adiabatik yang kemudian digunakan untuk membangun sebuah siklus mesin panas kuantum. \cite{Vinjanampathy2016Quantum}
\par Fraktal merupakan sebuah kajian matematis tentang geometri yang memiliki dimensi lebih besar dari dimensi topologisnya, umumnya bernilai pecahan. Dimensi tersebut disebut sebagai dimensi fraktal \cite{Mandelbrot1982Fractal}. Konsep dari fraktal ini diterapkan ke banyak aspek, salah satunya adalah fisika kuantum. Penerapan awal fraktal terhadap kuantum dilakukan oleh Feynman melalui pendekatan \textit{path integral}-nya terhadap persamaan Schrodinger. Yang mana merupakan penerapan gerakan Brownian terhadap fisika kuantum \cite{Feynman1948Quantum}. Pekerjaan ini dilanjut pada tahun 2000-an, di mana Nick Laskin menerbitkan tulisan yang secara eksplisit menerangkan penerapan kalkulus fraksional terhadap fisika kuantum \cite{Laskin2000Fractional}. Pendekatan yang dibuat Laskin ini merupakan generalisasi pekerjaan Feynman karena menggunakan Levy \textit{flight} daripada gerakan Brownian. Gerakan Brownian tersendiri merupakan gerakan dengan banyak langkah kecil yang memiliki arah yang acak untuk setiap langkahnya. Levy \textit{flight} adalah bentuk yang lebih umum dari gerakan Brownian di mana di antara langkah-langkah kecil diselingi beberapa lompatan besar. Kedua jenis gerakan ini termasuk ke dalam fraktal acak. \cite{Laskin2018Fractional}
\par Penerapan fraktal pada fisika kuantum ini tentu dapat diterapkan juga pada mesin panas kuantum menjadi mesin panas kuantum fraksional. Tentu sudah ada upaya serupa untuk menerapkan fraktal ke dalam mesin panas kuantum. Salah satu di antaranya adalah karya Ekrem Aydiner pada tahun 2021 yang berjudul, "Space-Fractional Quantum Heat Engine Based on Level Degeneracy." Jurnal tersebut membahas kerja serta efisiensi dari siklus Szillard kuantum. Pada jurnal tersebut ditemukan bahwa parameter fraksional memainkan peran yang penting dalam efisiensi dan kerja dari mesin yang ditinjau \cite{Aydiner2021Space}. Meski upaya yang sama telah dilakukan, masih banyak siklus yang bekum ditinjau. Maka dari itu, tulisan ini melakukan penerapan mesin panas kuantum fraksional yang serupa, namun untuk siklus yang berbeda. Selain itu, tinjauan juga tidak hanya dilakukan terhadap efisiensi, melainkan terhadap keluaran daya-nya juga. Untuk siklus yang ditinjau, didapatkan hasil yang sedikit berbeda di mana pengaruh parameter fraksional tidak terlalu besar terhadap efisiensi namun cukup signifikan untuk keluaran daya.
\par Terdapat 1 parameter penting lain yang juga ditinjau dalam tulisan ini. Parameter tersebut adalah parameter disipasi. Di sini dianggap kedua reservoir panas berukuran sangat besar. Sehingga dimungkinkan untuk terjadi kebocoroan atau disipasi panas di antara kedua reservoir. Hal ini mengacu pada metode pengerjaan mesin panas kuantum yang dilakukan oleh Jianhui Wang dan Jizhou He pada tahun 2012 dengan jurnal yang berjudul, "Optimization on a Three-Level Heat Engine Working with Two Nonnteracting Fermion in a One-Dimensional Box Trap." Pada jurnal tersebut, siklus yang ditinjau adalah siklus Carnot. Ditemukan bahwa parameter disipasi sangat mempengaruhi efisiensi dari siklus namun tidak mempengaruhi keluaran daya-nya \cite{Wang2012Optimization}. Ditemukan hasil yang sejalan pada tulisan ini, di mana efisiensi sangat dipengaruhi oleh disipasi panas sedangkan keluaran daya tidak terpengaruh sama sekali. 


 % Proses-proses ini tidak bisa diterapkan secara langsung terhadap sistem kuantum, harus dilakukan penyetaraan besaran terlebih dahulu. Dimulai dari isothermal yang berubah menjadi isoenergetik karena temperatur pada skala kuantum hanyalah energi rata-rata dari partikel. Untuk isovolume tidak terlalu berubah, namun jika yang ditinjau 1D maka volume diartikan sebagai lebar. Lalu isobarik, namanya tetap namun besaran yang ditinjau berbeda. Tekanan berarti gaya persatuan luas, namun luasan tidak memiliki makna berarti pada sistem kuantum 1 dimensi yang ada pada konteks tulisan ini. Sehingga besaran yang ditinjau adalah gaya yang bekerja pada tembok potensial. Lalu terakhir adalah proses adiabatik yang tidak jauh berbeda dengan versi klasiknya di mana tidak ada panas yang masuk dan keluar selama proses berlangsung. Namun proses ini bisa juga disebut sebagai isoentropis karena entropi selama proses ini konstan. \cite{Quan2008QHE}
% \par Sistem kuantum yang ditinjau pada tulisan ini adalah permasalahan potensial tak hingga 1 dimensi. Susunan dari permasalahan ini adalah sebuah partikel yang berada dalam sebuah tempat dengan potensial 0 berdimensi 1 dengan panjang \(L\) dan terjebak di antara 2 tembok potensial bernilai tak hingga. Potensial tak hingga ini berfungsi untuk memastikan partikel benar-benar terjebak tanpa ada probabilitas untuk keluar sama sekali \cite{GusPur2016Kuantum}. Untuk melakukan penerapan siklus mesin panas pada sistem kuantum ini, salah satu tembok potensial dibiarkan bergerak dengan bebas agar volume dalam sumur bisa berubah-ubah \cite{Latifah2013QHE}. Siklus mesin panas yang diterapkan pada tulisan ini ada 3 yaitu siklus Carnot, Otto, dan Brayton. Dengan menggunakan pendekatan Levy \textit{flight} dengan indeksnya yang pada konteks ini disebut parameter fraksional, bentuk persamaan tingkatan energi eigen ditentukan. Dengan menggunakan persamaan itu, efisiensi mesin serta daya yang dikeluarkan mesin panas kuantum diinvestigasi untuk melihat pengaruh parameter fraksional terhadap kedua hal tersebut. 

\section{Rumusan Masalah}
Berdasarkan latar belakang yang telah dijelaskan sebelumnya, rumusan masalah yang dikaji dalam tulisan ini adalah bagaimana pengaruh parameter fraksional terhadap efisiensi serta daya yang dikeluarkan oleh mesin panas kuantum.

\section{Tujuan}
Dari rumusan masalah yang telah disebutkan, tujuan dari pembuatan tulisan ini antara lain,
\begin{enumerate}
	\item Memahami bagaimana menerapkan konsep kalkulus fraksional dalam kasus mesin panas kuantum.
	\item Melihat pengaruh dari parameter fraksional pada keluaran daya serta efisiensi pada mesin panas kuantum.
        \item Melihat pengaruh dari parameter disipasi terhadap keluaran daya serta efisiensi pada mesin panas kuantum.
        % \item Membandingkan hasil penelitian ini dengan penelitian yang ada sebelumnya untuk memeriksa keabsahan klaim tentang pengaruh parameter fraksional terhadap efisiensi dan daya mesin. 
\end{enumerate}

\section{Batasan Masalah}
Batasan masalah yang ditetapkan untuk tulisan ini antara lain adalah;
\begin{enumerate}
    \item Zat kerja yang ditinjau pada mesin panas ini adalah satu partikel non relativistik.
    \item Sistem kuantum yang ditinjau adalah permasalahan potensial tak hingga satu dimensi dengan dua keadaan; \(n = 1,2\). 
\end{enumerate}

\section{Manfaat Penelitian}
Tulisan ini diharapkan dapat membuat pembaca memahami konsep fraktal serta penerapannya dalam mesin panas kuantum. Serta, diharapkan bisa mempejelas pengaruh parameter fraksional terhadap efisiensi juga keluaran daya pada mesin panas kuantum. Terakhir, tulisan ini diharapkan menginspirasi peneliti lain untuk melakukan penerapan kalkulus fraksional lebih lanjut terhadap mesin panas kuantum ataupun bidang kuantum lainnya.

\newpage
\thispagestyle{empty}
\vspace*{\fill}
    \begin{center}
	Halaman ini sengaja dikosongkan
    \end{center}
\vspace*{\fill}
\pagebreak